\chapter{1958年北京卷}

\begin{problem}

解方程 \(x-5=\sqrt{x+7}\)

\end{problem}

\begin{solution}
解：由原方程得 \(x^{2}-10x+25=x+7\)，即

\[
x^{2}-11x+18=0
\]

所以 \((x-2)(x-9)=0\)，得 \(x_{1}=2\)，\(x_{2}=9\). 

经检验 \(x=2\) 是增根，\(x=9\) 是原方程的解.
\end{solution}

\begin{problem}

已知 \(\lg 2=0.3010\)，试利用对数性质判定 \(2^{7}\cdot8^{11}\cdot5^{10}\) 是几位数

\end{problem}

\begin{solution}
解：
\[
\lg\left(2^{7}\cdot8^{11}\cdot5^{10}\right)
=\lg\left(2^{7}\cdot2^{33}\cdot\frac{10^{10}}{2^{10}}\right)
=\lg\left(2^{30}\cdot10^{10}\right)
=30\lg 2+10
=9.030+10
=19.03
\]

因此 \(2^{7}\cdot8^{11}\cdot5^{10}\) 是 20 位数.
\end{solution}

\begin{problem}

由50个不同的元素中选47个，问有多少种选法

\end{problem}

\begin{solution}
解：
\[
\mathrm C_{50}^{47}=\mathrm C_{50}^{3}
=\frac{50\times49\times48}{3\times2\times1}
=19600
\]

故有 19600 种选法.
\end{solution}

\begin{problem}

求 \(\operatorname{tg}\left(\operatorname{arctg}\frac56-\operatorname{arctg}\frac38\right)\) 之值

\end{problem}

\begin{solution}
解：令 \(\operatorname{arctg}\frac56=\alpha\)，\(\operatorname{arctg}\frac38=\beta\)

因此 \(\operatorname{tg}\alpha=\frac56\)，\(\operatorname{tg}\beta=\frac38\)

\[
\operatorname{tg}(\alpha-\beta)
=\frac{\operatorname{tg}\alpha-\operatorname{tg}\beta}{1+\operatorname{tg}\alpha\cdot\operatorname{tg}\beta}
=\frac{\frac56-\frac38}{1+\frac56\cdot\frac38}
=\frac{22}{63}
\]

因此
\[
\operatorname{tg}\left(\operatorname{arctg}\frac56-\operatorname{arctg}\frac38\right)=\frac{22}{63}
\]

又解
\[
\text{原式}=
\frac{\operatorname{tg}\left(\operatorname{arctg}\frac56\right)-\operatorname{tg}\left(\operatorname{arctg}\frac38\right)}
{1+\operatorname{tg}\left(\operatorname{arctg}\frac56\right)\cdot\operatorname{tg}\left(\operatorname{arctg}\frac38\right)}
=
\frac{\frac56-\frac38}{1+\frac56\times\frac38}
=\frac{22}{63}
\]
\end{solution}

\begin{problem}

在 \(\triangle ABC\) 中，已知 \(\angle A=135^\circ\)，\(\angle B=15^\circ\)，\(C=1\)，问三角形的最大边长是多少

\end{problem}

\begin{solution}
解：\(\angle C=180^\circ-(135^\circ+15^\circ)=30^\circ\)

\[
\frac{BC}{\sin135^\circ}=\frac{1}{\sin30^\circ}
\]

\[
BC=\frac{\sin135^\circ}{\sin30^\circ}
=\frac{\sin45^\circ}{\sin30^\circ}
=\sqrt2
\]

故最大边长为 \(\sqrt2\).
\end{solution}

\begin{problem}

过球面上的两个不同点能作几个大圆，为什么

\end{problem}

\begin{solution}
答：如果这两个点的连线不过球心，只能作一个大圆，因为两个不同的点与圆心只决定一个平面，这个平面与球只交出一个大圆. 

如果球心在两点的连线上，那么过这条直线可作无数个平面，与球交出无数个大圆.
\end{solution}

\begin{problem}

自 \(30^\circ\) 的二面角的一个面上一点 \(A\) 作另一个面的垂线，垂足是 \(B\)，若已知 \(AB\) 的长是 \(a\)，求 \(A\) 点到二面角的棱的距离

\end{problem}

\begin{solution}
解：作 \(BC\perp MN\)，连接 \(AC\)，则 \(AC\perp MN\)（三垂线定理）

因此 \(\angle ACB\) 是二面角的平面角，\(\angle ACB=30^\circ\)

\bitmapfigure[max width=0.42\linewidth,max totalheight=4.8cm]{img/1958/beijing/q7_solution_fig1.png}

在直角 \(\triangle ABC\) 中，\(AB=a\)，\(\angle ACB=30^\circ\)，\(\angle ABC=90^\circ\)，所以 \(AC=2a\)

故 \(A\) 点到二面角的棱的距离为 \(2a\).
\end{solution}

\begin{problem}

已知二次函数 \(y=3+2x-x^{2}\)

（1）问 \(x\) 在什么范围内 \(y>0\) 又 \(x\) 在什么范围内 \(y<0\)

（2）求出这个函数图象的顶点坐标和对称轴方程

（3）画出这个函数的图象

\end{problem}

\begin{solution}
解：（1）由 \(3+2x-x^{2}>0\) 得 \((x+1)(x-3)<0\)，所以 \(-1<x<3\). 

由 \(3+2x-x^{2}<0\) 得 \((x+1)(x-3)>0\)，所以 \(x<-1\) 或 \(x>3\). 

故当 \(x\in(-1,3)\) 时 \(y>0\)，而 \(x\in(-\infty,-1)\cup(3,+\infty)\) 时 \(y<0\). 

（2）\(y=3+2x-x^{2}=-(x-1)^{2}+4\). 

顶点 \((1,4)\). 

对称轴方程 \(x=1\). 

（3）画出这个函数的图象

\bitmapfigure[max width=0.36\linewidth,max totalheight=4.8cm]{img/1958/beijing/q8_solution_fig1.png}
\end{solution}

\begin{problem}

设 \(\sin\alpha\) 和 \(\sin\beta\) 是方程 \(x^{2}-\left(\sqrt2\cos20^\circ\right)x+\left(\cos^{2}20^\circ-\frac12\right)=0\) 的两根，其中 \(\alpha\)，\(\beta\) 是锐角，求 \(\alpha\)，\(\beta\) 的度数

\end{problem}

\begin{solution}
解：
\[
x=\frac{\sqrt2\cos20^\circ\pm\sqrt{2(1-\cos^{2}20^\circ)}}{2}
=\frac{\sqrt2}{2}\left(\cos20^\circ\pm\sin20^\circ\right)
=\frac{\sqrt2}{2}\left(\sin70^\circ\pm\sin20^\circ\right)
\]

不妨设
\[
\sin\alpha=\frac{\sqrt2}{2}\left(\sin70^\circ+\sin20^\circ\right)
=\sqrt2\sin45^\circ\cdot\cos25^\circ
=\cos25^\circ
=\sin65^\circ
\]

所以 \(\alpha=65^\circ\)

\[
\sin\beta=\frac{\sqrt2}{2}\left(\sin70^\circ-\sin20^\circ\right)
=\sqrt2\cos45^\circ\cdot\sin25^\circ
=\sin25^\circ
\]

所以 \(\beta=25^\circ\)
\end{solution}

\begin{problem}

有甲、乙两容器，甲容器是直圆柱形，高2尺，底半径1尺，乙容器是直圆锥形（锥底向上）高2尺，底面半径 \(2\sqrt3\) 尺，若将甲容器灌满水，然后把甲容器的一部分水倒入乙容器使得甲乙两个容器的水面一样高，问这时水面的高度是多少

\end{problem}

\begin{solution}
解：设甲容器的容积为 \(V\)，甲倒入乙剩下来的体积为 \(V_{1}\)，倒入乙容器内水的体积为 \(V_{2}\)，并设这时水面高度为 \(h\)

\bitmapfigure[max width=0.62\linewidth,max totalheight=4.8cm]{img/1958/beijing/q10_solution_fig1.png}

乙容器内水面的半径 \(R\) 满足 \(\frac{R}{2\sqrt3}=\frac h2\)，所以 \(R=\sqrt3 h\). 

所以 \(V_{1}=\pi\cdot1^{2}\cdot h=\pi h\)，\(V_{2}=\frac13\pi(\sqrt3 h)^{2}\cdot h=\pi h^{3}\). 

又 \(V_{1}+V_{2}=V\)，即 \(\pi h+\pi h^{3}=\pi\cdot1^{2}\cdot2=2\pi\)，故 \(h^{3}+h-2=0\)

设 \(f(h)=h^{3}+h-2\)

因为 \(f(1)=0\)，所以 \(h=1\) 是方程 \(h^{3}+h-2=0\) 的一个根

\[
h^{3}+h-2=(h-1)(h^{2}+h+2)=0
\]

而方程 \(h^{2}+h+2=0\) 无实根

故 \(h^{3}+h-2=0\) 只有一个实数根 \(h=1\)

所以甲乙容器水面一样高时，这时水面高度是 1 尺.
\end{solution}

\begin{problem}

已知 \(\widehat{AB}\) 是 \(\odot O\) 上含 \(120^\circ\) 的弧，过 \(A\)、\(B\) 两点所引 \(\odot O\) 的切线交于 \(C\)，\(\odot M\) 是与 \(\widehat{AB}\)、\(AC\)、\(BC\) 都相切的圆，求证 \(\odot M\) 的周长等于 \(\odot O\) 的周长的 \(\frac13\)

\end{problem}

\begin{solution}
解：因为 \(\widehat{AB}\) 是 \(120^\circ\)，所以 \(\angle AOC=60^\circ\)

所以 \(\angle ACO=30^\circ\)

因此 \(OC=2R\)（在直角三角形中，斜边等于 \(30^\circ\) 角所对边长的 2 倍）

所以 \(CM=2R-(R+r)=R-r\)

\bitmapfigure[max width=0.45\linewidth,max totalheight=4.8cm]{img/1958/beijing/q11_solution_fig1.png}

\[
\triangle CNM\sim\triangle AOC
\]

\[
\frac{r}{R}=\frac{R-r}{2R}
\]

所以 \(r=\frac R3\)

设 \(\odot M\) 的周长为 \(C_{\odot M}\)，\(\odot O\) 的周长为 \(C_{\odot O}\). 

\[
C_{\odot M}=2\pi r=2\pi\cdot\frac R3=\frac13\left(2\pi R\right)=\frac13 C_{\odot O}
\]

因此 \(\odot M\) 的周长等于 \(\odot O\) 的周长的 \(\frac13\).
\end{solution}
